\section{总结与展望}
\label{sec:conclusion}

\subsection{工作总结}

本文研究正常模板可用条件下的PCB检测结果复核，并构建由候选生成、多模态验证器、复检门和接受门组成的后置系统。实验同时给出有效组件、无显著增量的设计和端到端能力边界。

系统合并基础检测候选与差异候选，使复核范围不受基础检测器输出限制。阶段A调参后，差异候选在IoU 0.3下的漏检恢复率达到0.625。Oracle-C的F1为0.9555，相对基础方法保留约0.044的纠错空间。双来源候选覆盖了一部分基础检测器未提出的位置，但候选覆盖不等同于端到端纠错收益。

系统使用配准质量、检测不确定性和候选来源进行分级复核。在5\%调用预算下，Logistic复检门覆盖46.49\%的可纠错错误，随机选择覆盖3.15\%，前者约为后者的14.8倍。$q_{\mathrm{reg}}$在选定的同方向平移序列中随扰动增强而下降，可作为配准可靠性指标。强扰动条件下的HCR为1.0，但分母只有1个被接受动作。冻结阈值0.5的配准质量规则没有拦截任何有害动作。因此，当前实验尚未验证该规则的保护作用，也不支持质量分在所有几何扰动间严格单调。

学习式接受门和硬安全规则约束验证器动作。完整系统的正确检测保留率为0.9998，734个正确检测中误删2个。双门控将HCR从无门控全量复核的0.9326降至0.4000，并把误删控制在2次，仍有有害修改发生。加入接受门后，各调用预算下的F1与B1无统计差异，保持在基线水平。该结果支持双门控大幅减少有害修改，不能解释为系统消除了动作风险。

显式外观差异图与拓扑图没有形成预期贡献。200候选主消融中，A4为0.7571，略高于A2的0.7553，但差异处于置信区间内，未证明拓扑图带来增量。A5和A6相对A0的置信区间均跨零，因此不显著。阶段B2的450次专门对照采用另一批150个候选，A2为0.8082，并高于全部含差异图的变体。两个实验口径不能混用。最终系统移除显式差异视觉输入，只保留拓扑统计作为门控特征。去锚定设计也未获支持，提供检测标签的对照更好。

\subsection{结论与能力边界}

双门控大幅减少了有害修改，但端到端净收益接近零。五折实验中，完整系统F1为$0.9117\pm0.0343$，基础方法为$0.9118\pm0.0344$。二者的配对bootstrap差值为$-0.000064$，95\%置信区间为$[-0.000273,+0.000094]$，区间跨零。完整系统mAP@0.5为0.8810，基础方法为0.8813。当前证据不支持“复核系统提升整体检测精度”的结论。

Oracle分析进一步定位了限制因素。Oracle-C的F1为0.9555，说明候选集合仍有纠错空间。即使接受门理想地选择当前验证器提出的动作，Oracle-A的F1也只有0.9143，仅比基础方法高0.0025。当前主要限制因素是验证器对工业细粒度缺陷的判断和动作提议质量，而非候选覆盖。实验支持复检门提高有限预算下的错误候选覆盖效率，也支持接受门减少有害修改。实验不支持整体精度提升，端到端净收益接近零。

\subsection{未来工作}

接受门需要匹配真实提议分布。现有接受门使用确定性证据代理产生的动作训练，实验已积累约7000次真实验证器调用及其缓存输出。后续可据此重建动作级训练集，并按模板组隔离训练与测试。新模型仍需以正确检测保留率和有害修正率作为硬约束。

跨后端验证需要至少一个独立的多模态验证器，并在相同候选、证据卡、提示和预算下进行配对比较。跨模型一致的失败可用于检查任务或证据设计。模型间差异可用于分析领域能力、校准和动作偏好。当前单一后端的结论不能泛化到全部多模态验证器。

跨数据集实验可在HRIPCB上按模板分组，避免同一模板跨训练集与测试集造成泄漏。评价需分别报告配准成功率、差异候选召回、门控覆盖率与最终检测指标。该实验用于检验当前DeepPCB结果能否迁移到旋转变化和更多模板条件。

端到端增益取决于工业细粒度缺陷的领域适配。候选级正常与缺陷图对、缺陷类别知识和真实纠错动作可用于监督微调或偏好优化，第二模型或外部视觉专家可提供对照反馈。领域适配后的验证器仍需经过独立接受门，模型能力不能替代系统安全约束。
